\section{Introduction}

Recent advances in 3D vision and robot learning have significantly expanded the capabilities of embodied intelligent systems. On one hand, neural rendering and 3D generation techniques enable the reconstruction of realistic virtual environments from images. On the other hand, learning-based manipulation policies allow robots to acquire complex skills directly from demonstrations. Together, these developments provide a foundation for building interactive virtual worlds and autonomous robotic systems that can perceive and act in complex environments.

This project consists of two complementary tasks. The first task focuses on 3D reconstruction and scene generation. Multiple objects and background components are reconstructed using a combination of multi-view reconstruction and image-guided 3D generation methods, and the resulting assets are integrated into a complete virtual scene. The second task investigates the cross-environment generalization ability of robotic manipulation policies. Specifically, we evaluate whether Action Chunking with Transformers (ACT) can successfully transfer to visually different environments under a zero-shot setting.

In Task 1, several 3D assets are reconstructed. For Object A and the background, we use COLMAP for camera pose estimation combined with 2D Gaussian Splatting (2DGS) to exploit multi-view geometric consistency. For Objects B and C, we leverage threestudio and Magic123 for image-guided 3D generation. We analyze training dynamics through loss curves and evaluate reconstruction quality using standard metrics including SSIM, PSNR, LPIPS, and training time. Finally, the reconstructed assets are exported as textured meshes and integrated into a complete scene in Blender.

In Task 2, we study the zero-shot transfer capability of ACT on the CALVIN benchmark under substantial visual distribution shifts. Policies are trained using demonstrations collected from one or multiple environments and subsequently evaluated on a previously unseen environment. We assess policy performance using both offline action prediction metrics and online rollout success rates, and further investigate the influence of multi-environment training and action chunking on cross-environment generalization.
